%===============================================================
%  Rapport de stage -- ETNIC 2026
%  Patrick PREDA -- Master cybersecurite ULB -- STAG-Y-006
%===============================================================
\documentclass{rapport}

\usepackage{lipsum}
\usepackage[utf8]{inputenc}
\usepackage[T1]{fontenc}
\usepackage[french]{babel}
\usepackage{color}
\usepackage{xcolor}
\usepackage{graphicx}
\usepackage{booktabs}
\usepackage{hyperref}
\usepackage{listings}
\usepackage{minted}
\usepackage{tcolorbox}
\usepackage{pgfgantt}
\usepackage{geometry}
\usepackage{enumitem}
\usepackage{tabularx}
\usepackage{array}
\usepackage{multirow}
\usepackage{float}

\tcbuselibrary{skins, breakable}

% Couleurs ETNIC / ULB
\definecolor{ulbblue}{RGB}{0,68,139}
\definecolor{etnicgreen}{RGB}{0,150,100}
\definecolor{codebg}{RGB}{245,245,245}
\definecolor{alertred}{RGB}{200,50,50}

% Boite technique
\newtcolorbox{boitetech}[2][]{
  colback=codebg,
  colframe=ulbblue,
  fonttitle=\bfseries,
  title={#2},
  breakable,
  #1
}

% Config listings
\lstset{
  basicstyle=\ttfamily\small,
  backgroundcolor=\color{codebg},
  frame=single,
  breaklines=true,
  keywordstyle=\color{ulbblue}\bfseries,
  commentstyle=\color{gray}\itshape,
  stringstyle=\color{etnicgreen},
}

\title{Rapport de stage -- ETNIC}

\begin{document}

%===============================================================
%  PAGE DE GARDE
%===============================================================
\logo{logos/ulbLogo.png}
\titre{Rapport de stage\\[0.4em]
       \large Cybersecurite operationnelle a l'ETNIC\\[0.2em]
       \normalsize 2 mars 2026 -- 29 mai 2026}
\cours{STAG-Y-006}

\enseignant{
  \textbf{Maitre de stage :}\\ Laurent \textsc{Servais}

  \textbf{Promoteur ULB :} \\Jerome \textsc{Dossogne}
}

\eleves{
  Patrick \textsc{Preda}\\
  Noma : \textsc{615608}
}

\fairemarges
\fairepagedegarde
\tabledematieres

%===============================================================
%  REMERCIEMENTS
%===============================================================
\newpage
\section*{Remerciements}
\addcontentsline{toc}{section}{Remerciements}

Je tiens a remercier Laurent \textsc{Servais}, responsable du Centre de
Competences Securite d'ETNIC, pour son encadrement et la confiance accordee
tout au long de ce stage. Je remercie egalement Tom \textsc{Pirenne} et
l'ensemble des equipes P3SI, IAM et TAR pour leur disponibilite et leur
partage de connaissances. Mes remerciements vont aussi a Jerome
\textsc{Dossogne} de l'ULB pour son suivi academique.

%===============================================================
%  LISTE DES ABREVIATIONS
%===============================================================
\newpage
\section*{Liste des abreviations}
\addcontentsline{toc}{section}{Liste des abreviations}

\begin{tabularx}{\linewidth}{lX}
\toprule
\textbf{Abreviation} & \textbf{Definition} \\
\midrule
CERT-FR  & Centre gouvernemental francais de veille, d'alerte et de reponse aux attaques informatiques \\
CMDB     & Configuration Management Database \\
EDR      & Endpoint Detection et Response \\
EPP      & Endpoint Protection Platform \\
ETNIC    & Entreprise des Technologies Nouvelles de l'Information et de la Communication \\
FWB      & Federation Wallonie-Bruxelles \\
IAM      & Identity et Access Management \\
IOC      & Indicator of Compromise \\
MITRE    & MITRE ATT\&CK : referentiel des techniques d'attaques adversariales \\
NPM      & Nginx Proxy Manager \\
P3SI     & Produits et Services de Securite des Systemes d'Information \\
PNDS     & PowerShell Network Diagnostic Script \\
RSS      & Really Simple Syndication \\
SCADA    & Supervisory Control and Data Acquisition \\
SIEM     & Security Information and Event Management \\
SNOW     & ServiceNow \\
SOC      & Security Operations Center \\
SOAR     & Security Orchestration, Automation and Response \\
TAR      & Tracabilite, Audit et Reporting \\
TLS      & Transport Layer Security \\
VPN      & Virtual Private Network \\
\bottomrule
\end{tabularx}

%===============================================================
%  1. INTRODUCTION
%===============================================================
\newpage
\section{Introduction}

Ce rapport presente le deroulement de mon stage de fin d'etudes effectue au
sein de l'\textbf{ETNIC} (Entreprise des Technologies Nouvelles de
l'Information et de la Communication) du \textbf{2 mars au 29 mai 2026},
dans le cadre du Master en cybersecurite a finalite \textit{Conception et
Analyse de Systemes} de l'Universite libre de Bruxelles (ULB).

L'ETNIC est l'operateur informatique de la Federation Wallonie-Bruxelles.
Son Centre de Competences Securite assure la gestion de la securite
operationnelle de l'ensemble du systeme d'information de la FWB, couvrant
des milliers d'assets, d'utilisateurs et de services critiques.

Ce stage visait a m'immerger dans un SOC operationnel reel et a contribuer
concretement a l'amelioration de la posture de securite de l'organisation.
J'ai eu l'opportunite de travailler sur cinq missions distinctes : la veille
securite, l'automatisation d'inventaires, la sensibilisation au phishing, le
diagnostic reseau et le monitoring de flux de threat intelligence.

Le rapport est structure comme suit : apres une presentation de l'organisation
et des objectifs du stage, une timeline recapitule le deroulement
chronologique. Chaque mission est ensuite decrite techniquement, suivie d'un
bilan des competences acquises.

%===============================================================
%  2. PRESENTATION DE L'ETNIC
%===============================================================
\section{Presentation de l'ETNIC}

\subsection{Contexte organisationnel}

L'ETNIC est l'operateur ICT de la Federation Wallonie-Bruxelles, cree pour
mutualiser les ressources informatiques des organismes publics francophones
belges. Son siege est situe Boulevard du Roi Albert II, 37 a 1030 Bruxelles.

\begin{boitetech}{Chiffres cles}
\begin{itemize}[noitemsep]
  \item Entite publique au service de la FWB et de ses organismes
  \item Gestion centralisee de l'infrastructure IT et de la cybersecurite
  \item Plusieurs milliers d'assets geres dans le CMDB ServiceNow
  \item Infrastructure hybride : on-premise et cloud
\end{itemize}
\end{boitetech}

\subsection{Centre de Competences Securite}

Le stage s'est deroule au sein du \textbf{Centre de Competences Securite
(CCS)}, qui regroupe trois equipes operationnelles :

\begin{itemize}
  \item \textbf{P3SI} : Produits et Services de Securite des Systemes
        d'Information, gestion des outils de securite (SIEM, EDR, scanner
        de vulnerabilites, etc.) ;
  \item \textbf{IAM} : Identity et Access Management, gestion des identites,
        des acces et des politiques d'authentification ;
  \item \textbf{TAR} : Tracabilite, Audit et Reporting, production des rapports
        de securite et suivi des indicateurs.
\end{itemize}

J'ai principalement travaille avec l'equipe \textbf{P3SI}, sous la
responsabilite de Laurent \textsc{Servais}.

%===============================================================
%  3. OBJECTIFS DU STAGE
%===============================================================
\section{Objectifs du stage}

\subsection{Objectifs officiels (convention STAG-Y-006)}

La convention de stage definit les missions suivantes :

\begin{itemize}
  \item Contribuer aux activites operationnelles du service cybersecurite ;
  \item Ameliorer les capacites de detection et de reponse aux evenements de
        securite ;
  \item Participer a l'optimisation de la gestion des vulnerabilites ;
  \item Participer a l'optimisation des services EPP/EDR et WASC ;
  \item Integrer une approche orientee menace dans les processus de securite ;
  \item Developper l'usage des API de solutions de securite et le scripting
        d'automatisation.
\end{itemize}

\subsection{Objectifs personnels}

Au-dela des objectifs contractuels, je souhaitais personnellement :

\begin{itemize}
  \item Acquerir une experience concrete en SOC dans un environnement de
        production ;
  \item Developper des competences en developpement d'outils de securite
        (Python, TypeScript) ;
  \item Comprendre les processus operationnels d'un CSIRT/SOC francophone ;
  \item Confronter les concepts academiques (MITRE ATT\&CK, SIEM, SOAR) a
        la realite terrain.
\end{itemize}

%===============================================================
%  4. DEROULEMENT DU STAGE
%===============================================================
\section{Deroulement du stage}

\subsection{Vue chronologique}

Le stage s'est etale sur \textbf{12 semaines effectives} (13 semaines
calendaires dont une semaine de suspension du 13 au 17 avril 2026).

\begin{figure}[H]
\centering
\begin{ganttchart}[
    hgrid,
    vgrid={dotted},
    x unit=0.72cm,
    y unit title=0.6cm,
    y unit chart=0.55cm,
    title/.style={fill=ulbblue!20, draw=ulbblue},
    title label font=\small\bfseries,
    bar/.style={fill=ulbblue!60, rounded corners=2pt},
    group/.style={fill=ulbblue, draw=ulbblue},
  ]{1}{13}
  \gantttitle{Stage ETNIC, Mars--Mai 2026}{13} \\
  \gantttitlelist{1,...,13}{1} \\

  \ganttbar{Onboarding et Lecture}{1}{2} \\
  \ganttbar{CertFR Monitor}{3}{3} \\
  \ganttbar{GoPhish, phishing interne}{4}{6} \\
  \ganttbar{PNDS, diagnostic reseau}{6}{8} \\
  \ganttbar[bar/.style={fill=alertred!40, draw=alertred,
    rounded corners=2pt}]{Suspension}{7}{7} \\
  \ganttbar{RapidETNIC}{9}{12} \\
  \ganttbar{Redaction rapport}{12}{13}
\end{ganttchart}
\caption{Diagramme de Gantt -- deroulement du stage}
\label{fig:gantt}
\end{figure}

\subsection{Phases}

\begin{description}[leftmargin=2cm, style=nextline]
  \item[Semaines 1--2 \small(2--13 mars)]
    \textbf{Onboarding et phase de lecture.}
    Prise en main de l'environnement ETNIC, des outils et des processus
    internes. Lecture approfondie de documents de reference sur le SIEM,
    le SOAR et le SOC.

  \item[Semaine 3 \small(16--20 mars)]
    \textbf{Moniteur CERT-FR.}
    Developpement d'un dashboard de veille des publications CERT-FR avec
    alertes email, backend Node.js/TypeScript et frontend Next.js.

  \item[Semaines 4--6 \small(23 mars -- 10 avr.)]
    \textbf{Campagne de sensibilisation au phishing (GoPhish).}
    Deploiement d'une infrastructure GoPhish complete avec reverse proxy
    Cloudflare/NPM.

  \item[Semaines 6 et 8 \small(6--10 avr. et 20--24 avr.)]
    \textbf{Script de diagnostic reseau (PNDS).}
    Developpement d'un script PowerShell pour le Service Desk de l'ETNIC,
    interrompu par la semaine de suspension et finalise a la reprise.

  \item[Semaine 7 \small(13--17 avr.)]
    \textit{Suspension du stage (accord avec le maitre de stage).}

  \item[Semaines 9--12 \small(27 avr. -- 22 mai)]
    \textbf{RapidETNIC.}
    Developpement d'un outil de croisement Rapid7 InsightVM et ServiceNow
    CMDB.

  \item[Semaines 12--13 \small(18--29 mai)]
    \textbf{Redaction du rapport de stage.}
\end{description}

%===============================================================
%  5. PROJETS REALISES
%===============================================================
\section{Projets realises}

%---------------------------------------------------------------
\subsection{Onboarding et phase de lecture}
%---------------------------------------------------------------

\subsubsection{Contexte}

Les deux premieres semaines ont ete consacrees a la prise en main de
l'environnement de travail et a une phase de lecture structuree. L'objectif
etait d'acquerir une base solide sur les technologies et approches en usage
dans l'equipe avant de passer aux missions operationnelles.

\subsubsection{Documents etudies}

\begin{table}[H]
\centering
\begin{tabularx}{\linewidth}{p{6.5cm}X}
\toprule
\textbf{Document} & \textbf{Apport principal} \\
\midrule
\textit{SOAR: How to Identify and Implement Security Automation Use Cases}
  & Methodologie de selection et d'implementation de cas d'usage SOAR ;
    criteres de priorisation des automatisations \\
\addlinespace
\textit{A Guidance Framework for Architecting and Deploying a Modern SIEM Solution}
  & Architecture cible d'un SIEM moderne ; sources de donnees, regles de
    correlation, integration dans le SOC \\
\addlinespace
\textit{How to Create and Maintain Security Monitoring Use Cases for Your SIEM}
  & Cycle de vie des use cases de detection ; formalisation, test et
    maintenance \\
\addlinespace
\textit{Cyberun \#44}
  & Panorama des menaces et actualites cybersecurite francophones \\
\bottomrule
\end{tabularx}
\caption{Documents de reference etudies en phase d'onboarding}
\end{table}

\subsubsection{Points cles retenus}

\begin{boitetech}{Synthese SIEM et SOAR}
\begin{itemize}[noitemsep]
  \item Un SIEM efficace repose sur la qualite des sources de donnees
        ingeres et la pertinence des regles de correlation, pas sur leur
        quantite.
  \item Le SOAR automatise les taches repetitives (enrichissement IOC,
        blocage IP, notification) pour liberer les analystes sur les
        incidents complexes.
  \item La maintenance des use cases est aussi critique que leur creation :
        un use case non revu genere du bruit ou des faux negatifs.
  \item L'approche \textit{threat-informed defense} via MITRE ATT\&CK guide
        le choix des detections prioritaires en fonction des TTPs
        adversariaux observes.
\end{itemize}
\end{boitetech}

%---------------------------------------------------------------
\subsection{Moniteur CERT-FR}
%---------------------------------------------------------------

\subsubsection{Contexte et objectifs}

Le CERT-FR publie quotidiennement des bulletins de securite critiques :
alertes, avis, IOC, bulletins SCADA, threat intelligence. L'ETNIC, comme
de nombreuses organisations francophones, suit ces publications pour
prioriser ses actions de remediation.

L'objectif etait de developper un outil centralise permettant de :
\begin{itemize}
  \item Agreger en temps reel tous les flux RSS du CERT-FR ;
  \item Alerter par email a chaque nouvelle publication ;
  \item Offrir un dashboard web filtrable pour la consultation.
\end{itemize}

\subsubsection{Architecture}

\begin{boitetech}{Stack technique}
\begin{tabular}{ll}
\textbf{Composant} & \textbf{Technologie} \\
Base de donnees    & PostgreSQL 16 \\
Backend            & Node.js / TypeScript / Express \\
Frontend           & Next.js 14 \\
Orchestration      & Docker Compose \\
Alerting email     & SMTP ou Resend API \\
\end{tabular}
\end{boitetech}

\begin{figure}[H]
\centering
\begin{verbatim}
CERT-FR RSS Feeds (7 flux)
         |
         v
  Backend Node.js/TS
  |- Cron job (toutes les 15 min)
  |- Dedoublonnage par GUID
  |- Stockage PostgreSQL
  `- Alertes email (SMTP / Resend)
         |
         v
  Frontend Next.js
  `- Dashboard : items, filtres, pagination, logs
\end{verbatim}
\caption{Architecture de CertFR Monitor}
\end{figure}

\subsubsection{Flux surveilles}

\begin{table}[H]
\centering
\begin{tabular}{llp{7cm}}
\toprule
\textbf{Type} & \textbf{Label} & \textbf{Contenu} \\
\midrule
\texttt{alerte}    & Alertes de securite  & Vulnerabilites critiques actives \\
\texttt{avis}      & Avis de securite     & Bulletins de mise a jour \\
\texttt{cti}       & Menaces et incidents & Rapports de threat intelligence \\
\texttt{ioc}       & IOC                  & Indicateurs de compromission \\
\texttt{dur}       & Durcissement         & Recommandations de configuration \\
\texttt{scada}     & SCADA                & Bulletins systemes industriels \\
\texttt{actualite} & Actualites           & Bulletins d'information generaux \\
\bottomrule
\end{tabular}
\caption{Flux RSS CERT-FR surveilles}
\end{table}

\subsubsection{API REST exposee}

\begin{lstlisting}[language=bash, caption=Endpoints API CertFR Monitor]
GET  /api/items   # Liste paginee (filtres: feedType, limit, offset)
GET  /api/stats   # Compteurs par flux
GET  /api/logs    # Historique des fetches
GET  /api/feeds   # Definitions des flux
POST /api/fetch   # Declenchement manuel d'un fetch
GET  /api/health  # Health check
\end{lstlisting}

%---------------------------------------------------------------
\subsection{Campagne de sensibilisation au phishing (GoPhish)}
%---------------------------------------------------------------

\subsubsection{Contexte et objectifs}

La sensibilisation des utilisateurs au phishing est un axe majeur de la
strategie de securite de l'ETNIC. L'objectif de cette mission etait de
deployer une infrastructure de campagne de phishing simule en interne,
permettant de mesurer le taux de clics sur des emails frauduleux et d'en
tirer des indicateurs pour orienter les formations.

\subsubsection{Architecture deployee}

\begin{boitetech}{Stack technique}
\begin{itemize}[noitemsep]
  \item \textbf{GoPhish} : plateforme open-source de gestion de campagnes
        phishing
  \item \textbf{Docker Compose} : conteneurisation de l'application
  \item \textbf{Nginx Proxy Manager (NPM)} : reverse proxy TLS partage
  \item \textbf{Cloudflare} : DNS et protection du domaine
        \texttt{preda.be}
\end{itemize}
\end{boitetech}

\begin{figure}[H]
\centering
\begin{verbatim}
Cloudflare (orange cloud)
       |
       v
Nginx Proxy Manager (npm-network)
  |- gophish.preda.be  -> gophish:3333  (admin UI)
  `- phish.preda.be    -> gophish:8080  (landing pages)
\end{verbatim}
\caption{Architecture de l'infrastructure GoPhish}
\end{figure}

L'interface d'administration est exposee sur \texttt{gophish.preda.be}
(acces restreint par liste d'acces NPM), tandis que les pages de phishing
simule sont servies sur \texttt{phish.preda.be} avec certificat TLS
Cloudflare origin.

\subsubsection{Resultats}

La mission s'est limitee a la mise en place et la validation technique de
l'infrastructure. Aucune campagne n'a ete lancee contre des utilisateurs
reels durant le stage : les contraintes organisationnelles (validation RH,
communication prealable, accord de la direction) necessaires a toute
campagne de sensibilisation interne n'ont pas pu etre reunies dans le delai
imparti.

\begin{itemize}
  \item Infrastructure complete deployee, testee et validee en environnement
        isole ;
  \item Scenarios de phishing simule crees et verifies (envoi test vers
        adresses internes de controle) ;
  \item TLS fonctionnel sur les deux domaines via Cloudflare origin cert ;
  \item Documentation technique remise a l'equipe P3SI pour permettre le
        lancement d'une campagne reelle lors d'une prochaine initiative de
        sensibilisation.
\end{itemize}

%---------------------------------------------------------------
\subsection{Script de diagnostic reseau (PNDS)}
%---------------------------------------------------------------

\subsubsection{Contexte et objectifs}

Le Service Desk de l'ETNIC traite regulierement des tickets lies a des
problemes reseau (DNS, VPN, connectivite). Dans la quasi-totalite des cas,
les memes informations de base manquent : configuration IP, etat des
services VPN/EDR, resolutions DNS. Le technicien doit relancer l'utilisateur
pour obtenir ces donnees, allongeant inutilement les delais de resolution.

L'objectif etait de developper un script PowerShell autonome, executable
sans droits administrateur, qui collecte automatiquement toutes ces
informations et genere un rapport pret a etre joint au ticket.

\subsubsection{Fonctionnalites}

\begin{table}[H]
\centering
\begin{tabularx}{\linewidth}{clX}
\toprule
\textbf{\#} & \textbf{Section} & \textbf{Donnees collectees} \\
\midrule
1  & Configuration IP    & \texttt{ipconfig /all} complet \\
2  & DNS systeme         & Resolution de tous les hotes internes et externes \\
3  & DNS internes forces & Resolutions sur chacun des 4 DNS internes ETNIC \\
4  & Ping                & Cibles internes (IP et noms) et externes \\
5  & Traceroute          & \texttt{tracert} vers ressources internes et externes \\
6  & Services            & F5 DNS Relay, Cisco Umbrella, AnyConnect VPN, BIG-IP \\
7  & Detection VPN       & Interfaces reseau et processus VPN actifs \\
8  & Table de routage    & \texttt{route print} complet \\
9  & Test HTTPS          & Code retour HTTP et URL finale \\
10 & Cache DNS           & \texttt{Get-DnsClientCache} \\
\bottomrule
\end{tabularx}
\caption{Sections du rapport genere par PNDS}
\end{table}

\subsubsection{Deploiement}

Le script est concu pour un double usage :
\begin{itemize}
  \item \textbf{Manuel} : execution directe via
        \texttt{powershell.exe -ExecutionPolicy Bypass -File
        DiagnosticReseau.ps1} ;
  \item \textbf{GPO} : deploiement de masse via Group Policy avec raccourci
        Bureau, permettant a n'importe quel utilisateur de lancer le
        diagnostic en un clic.
\end{itemize}

Le fichier de sortie est automatiquement horodate et depose sur le Bureau :
\texttt{Diagnostic\_Reseau\_YYYYMMDD\_HHmmss.txt}.

%---------------------------------------------------------------
\subsection{Croisement Rapid7 et ServiceNow (RapidETNIC)}
%---------------------------------------------------------------

\subsubsection{Contexte et objectifs}

L'ETNIC dispose de deux sources de verite pour son parc de serveurs :

\begin{itemize}
  \item \textbf{Rapid7 InsightVM} : scanner de vulnerabilites qui maintient
        un inventaire des assets effectivement scannes sur le reseau ;
  \item \textbf{ServiceNow CMDB} : base de configuration officielle listant
        tous les serveurs et leur statut de cycle de vie (\textit{Installe},
        \textit{Absent}, \textit{Mis hors service}, etc.).
\end{itemize}

Ces deux sources divergent regulierement : des serveurs decommissionnes
dans SNOW restent scannes dans Rapid7, tandis que des serveurs actifs dans
SNOW sont absents des scans. L'objectif de RapidETNIC est d'automatiser la
detection et le traitement de ces ecarts.

\subsubsection{Fonctionnement}

\begin{figure}[H]
\centering
\begin{verbatim}
fetch_r7.py           fetch_snow.py
     |                      |
     v                      v
 results/r7/          results/snow/
      \                    /
       \______fetch.py____/
                |
      Algorithme de matching
      (hostname primaire, IP publique fallback)
                |
       +--------+--------+
       |        |        |
    decomm   absent   matched
    in_r7    in_r7   outdated
       |
  API on-prem InsightVM
  - Tag "deprecated"
  - Suppression agent
\end{verbatim}
\caption{Pipeline RapidETNIC}
\end{figure}

\subsubsection{Algorithme de matching}

Le matching entre les deux sources repose sur le \textbf{hostname normalise}
comme cle primaire. Une specificite importante du CMDB ETNIC est que le
champ \texttt{name} de \texttt{cmdb\_ci\_server} contient le modele
materiel (ex. \texttt{POWEREDGE R730}), et non le hostname. Ce dernier se
trouve dans le champ \texttt{asset\_tag}, apres le dernier caractere
\texttt{/} :

\begin{minted}[bgcolor=codebg, fontsize=\small]{text}
SERVEUR_002046/SETNIC04J  ->  hostname = SETNIC04J
/SESX36                   ->  hostname = SESX36
SERVEUR_002424            ->  pas de hostname -> non-matche
\end{minted}

Un \textbf{fallback IP} est disponible mais uniquement pour les IPs
publiques : les plages privees (\texttt{10.x}, \texttt{172.16-31.x},
\texttt{192.168.x}) sont exclues pour eviter les faux positifs lies aux
collisions d'adresses entre VLANs.

\subsubsection{Categories de resultats}

\begin{table}[H]
\centering
\begin{tabularx}{\linewidth}{lXl}
\toprule
\textbf{Categorie} & \textbf{Signification} & \textbf{Action} \\
\midrule
\texttt{decommissioned\_in\_r7} & SNOW = Mis hors service, encore scanne
  & Tag + suppression agent \\
\texttt{installed\_not\_in\_r7} & SNOW = Installe, absent de Rapid7
  & Alerte \\
\texttt{matched\_outdated}      & Dans les deux, OS EOL ou scan obsolete
  & Alerte \\
\texttt{matched\_current}       & Dans les deux, a jour
  & OK \\
\texttt{in\_r7\_only}           & Scanne, aucun CI dans SNOW
  & A investiguer \\
\bottomrule
\end{tabularx}
\caption{Categories de resultats de RapidETNIC}
\end{table}

\subsubsection{Actions automatisees sur l'API InsightVM on-premise}

Pour chaque asset de la categorie \texttt{decommissioned\_in\_r7}, le
script effectue automatiquement deux actions via l'API REST de la console
InsightVM on-premise :

\begin{enumerate}
  \item \textbf{Application du tag \texttt{deprecated}} : creation du tag
        custom si absent, puis assignation via
        \texttt{PUT /api/3/assets/\{id\}/tags/\{tagId\}} ;
  \item \textbf{Suppression de l'agent scanner} : identification de l'agent
        par GUID (\texttt{source == "R7 Agent"}) dans le Cloud API,
        correspondance avec l'agent on-prem par la cle, puis
        \texttt{DELETE /api/3/agents/\{agentId\}}.
\end{enumerate}

L'asset lui-meme n'est pas supprime : les scans reseau continuent de le
detecter. Seul le lien d'agent est retire.

\begin{minted}[bgcolor=codebg, fontsize=\small]{python}
# Extrait : fallback IP uniquement si IP publique
r7 = r7_by_host.get(hostname)
if r7 is None and ip and not is_private_ip(ip):
    r7 = r7_by_ip.get(ip)
\end{minted}

\subsubsection{Resultats}

\begin{boitetech}{Metriques -- execution du 26 mai 2026}
\begin{itemize}[noitemsep]
  \item \textbf{Assets Rapid7 analyses :} 3 156
  \item \textbf{CIs ServiceNow traites :} 2 345
  \item \textbf{Assets decommissionnes encore scannes :} 6
        (tag \texttt{deprecated} applique + agent supprime)
  \item \textbf{Assets installes absents de Rapid7 :} 207
  \item \textbf{Assets matches a jour :} 1 249
  \item \textbf{Assets matches mais outdated (OS EOL ou scan > 90j) :} 52
  \item \textbf{Assets dans R7 sans CI SNOW :} 272
  \item \textbf{OS EOL detectes :} 169 (5,4\% du parc)
\end{itemize}
\end{boitetech}

%===============================================================
%  6. COMPETENCES ACQUISES ET BILAN
%===============================================================
\section{Competences acquises et bilan}

\subsection{Competences techniques}

\begin{table}[H]
\centering
\begin{tabularx}{\linewidth}{lXl}
\toprule
\textbf{Domaine} & \textbf{Competence} & \textbf{Niveau} \\
\midrule
SIEM/SOAR    & Architecture, use cases, automatisation
  & \textcolor{etnicgreen}{\textbf{+}} \\
API Security & Rapid7 InsightVM REST, ServiceNow REST
  & \textcolor{etnicgreen}{\textbf{++}} \\
Python       & Scripting, requetes API, traitement JSON
  & \textcolor{etnicgreen}{\textbf{++}} \\
PowerShell   & Scripting Windows, deploiement GPO
  & \textcolor{etnicgreen}{\textbf{+}} \\
Docker       & Compose, reseaux, reverse proxy
  & \textcolor{etnicgreen}{\textbf{++}} \\
Node.js/TS   & Express, cron jobs, TypeScript
  & \textcolor{etnicgreen}{\textbf{+}} \\
Next.js      & Dashboard React, SSR
  & \textcolor{etnicgreen}{\textbf{+}} \\
PostgreSQL   & Modelisation, dedoublonnage, requetes
  & \textcolor{etnicgreen}{\textbf{+}} \\
Phishing     & GoPhish, infrastructure de campagne
  & \textcolor{etnicgreen}{\textbf{++}} \\
Vuln. mgmt   & Rapid7 InsightVM, CMDB, cycle de vie assets
  & \textcolor{etnicgreen}{\textbf{++}} \\
\bottomrule
\end{tabularx}
\caption{Competences techniques acquises (+ : consolide, ++ : maitrise)}
\end{table}

\subsection{Competences transversales}

\begin{itemize}
  \item \textbf{Travail en autonomie} : chaque mission a ete confiee avec
        un niveau d'autonomie croissant, des specifications initiales a
        la livraison ;
  \item \textbf{Documentation technique} : redaction de READMEs, de
        rapports d'execution et de procedures utilisables par l'equipe ;
  \item \textbf{Communication} : restitution reguliere des avancees a
        Laurent Servais et presentation des outils aux equipes ;
  \item \textbf{Gestion de version} : utilisation systematique de Git
        (GitHub et GitLab interne ETNIC) pour tous les projets.
\end{itemize}

\subsection{Bilan personnel}

Ce stage a confirme mon interet pour la cybersecurite operationnelle, en
particulier l'intersection entre le developpement d'outils et les
operations de securite. J'ai pu constater que dans un SOC reel, la valeur
ajoutee ne reside pas uniquement dans la detection, mais egalement dans
l'automatisation des taches repetitives et la qualite des donnees
d'inventaire.

% [TODO: ajouter reflexions personnelles]

%===============================================================
%  7. CONCLUSION
%===============================================================
\section{Conclusion}

Ce stage de douze semaines a l'ETNIC m'a permis de contribuer concretement
a l'amelioration de la posture de securite d'un operateur ICT critique. A
travers cinq missions de natures complementaires (sensibilisation, outillage,
veille et gestion des vulnerabilites) j'ai pu mettre en pratique l'ensemble
des competences acquises durant mon Master en cybersecurite.

Les outils developpes (PNDS, CertFR Monitor, RapidETNIC) ont ete livres
operationnels et documentes, et sont destines a etre maintenus et reutilises
par les equipes P3SI.

Ce stage constitue une etape fondamentale dans ma formation : il m'a permis
de passer du modele theorique a la realite operationnelle d'un SOC, et de
mesurer l'importance de l'automatisation, de la rigueur documentaire et de
la collaboration dans un environnement de cybersecurite professionnel.

%===============================================================
%  BIBLIOGRAPHIE
%===============================================================
\newpage
\section*{Bibliographie}
\addcontentsline{toc}{section}{Bibliographie}

\begin{enumerate}
  \item SANS Institute. \textit{SOAR: How to Identify and Implement Security
        Automation Use Cases}. SANS Reading Room, 2023.
  \item Gartner. \textit{A Guidance Framework for Architecting and Deploying
        a Modern SIEM Solution}. Gartner Research, 2023.
  \item IBM Security. \textit{How to Create and Maintain Security Monitoring
        Use Cases for Your SIEM}. IBM, 2022.
  \item Cyberun. \textit{Cyberun \#44}. Cyberun Magazine, 2026.
  \item CERT-FR. \textit{Publications de securite}.
        \url{https://www.cert.ssi.gouv.fr/}. Consulte en avril--mai 2026.
  \item Rapid7. \textit{InsightVM REST API Documentation}.
        \url{https://help.rapid7.com/insightvm/en-us/api/}.
        Consulte en avril--mai 2026.
  \item ServiceNow. \textit{REST API Reference -- Table API}.
        \url{https://docs.servicenow.com}. Consulte en mars--mai 2026.
  \item MITRE ATT\&CK. \textit{Enterprise Matrix}.
        \url{https://attack.mitre.org/}. Consulte en 2026.
\end{enumerate}

%===============================================================
%  ANNEXES
%===============================================================
\newpage
\appendix
\section{Convention de stage}

La convention de stage originale signee entre l'ULB, l'ETNIC et le
stagiaire (STAG-Y-006) est disponible en document separe.

% [TODO: inserer scan ou referencer fichier externe]

\section{Extraits de code}

\subsection{RapidETNIC -- is\_private\_ip}

\begin{minted}[bgcolor=codebg, fontsize=\small]{python}
import ipaddress

_PRIVATE_NETS = [
    ipaddress.ip_network("10.0.0.0/8"),
    ipaddress.ip_network("172.16.0.0/12"),
    ipaddress.ip_network("192.168.0.0/16"),
    ipaddress.ip_network("127.0.0.0/8"),
    ipaddress.ip_network("169.254.0.0/16"),
]

def is_private_ip(ip: str) -> bool:
    try:
        addr = ipaddress.ip_address(ip)
        return any(addr in net for net in _PRIVATE_NETS)
    except ValueError:
        return True  # unparseable -> traiter comme prive
\end{minted}

\subsection{CertFR Monitor -- cron fetch (TypeScript)}

\begin{minted}[bgcolor=codebg, fontsize=\small]{typescript}
// Backend : polling RSS toutes les 15 minutes
cron.schedule(`*/${FETCH_INTERVAL_MINUTES} * * * *`, async () => {
  for (const feed of FEEDS) {
    const items = await fetchRSS(feed.url);
    for (const item of items) {
      await db.upsert(item, { onConflict: 'guid' });
      if (item.isNew) await sendAlert(item, feed);
    }
  }
});
\end{minted}

\end{document}
